%%%
%%% Radical "⾇"
%%%
\section*{Radical 136: ``⾇''}\addcontentsline{toc}{section}{Radical 136: ⾇}\addcontentsline{loh}{figure}{\#\#\#\# 136: ⾇}

%%%%%%%%%% 舜 %%%%%%%%%%
\subsection*{舜}\addcontentsline{loh}{figure}{舜}

\begin{Entry}{舜}{12}{⾇}
  \begin{Phonetics}{舜}{shun4}
    \definition*{s.}{Shun, o nome de um monarca lendário da China antiga | Sobrenome: Shun}
  \end{Phonetics}
\end{Entry}

%%%%%%%%%% 舞 %%%%%%%%%%
\subsection*{舞}\addcontentsline{loh}{figure}{舞}

\begin{Entry}{舞}{14}{⾇}
  \begin{Phonetics}{舞}{wu3}[][HSK 5]
    \definition[支,段,个]{s.}{dança | palco; domínio das atividades sociais}
    \definition{v.}{mover"-se como numa dança | dançar com algo nas mãos; brincar com | florescer; empunhar; brandir | esvoaçar | fazer malabarismos; brincar com}
  \end{Phonetics}
\end{Entry}

\begin{Entry}{舞厅}{14,4}{⾇,⼚}
  \begin{Phonetics}{舞厅}{wu3ting1}[][HSK 7-9]
    \definition[家,间]{s.}{salão de dança; salão de baile}
  \end{Phonetics}
\end{Entry}

\begin{Entry}{舞厅舞}{14,4,14}{⾇,⼚,⾇}
  \begin{Phonetics}{舞厅舞}{wu3ting1wu3}
    \definition{s.}{dança de salão}
  \end{Phonetics}
\end{Entry}

\begin{Entry}{舞台}{14,5}{⾇,⼝}
  \begin{Phonetics}{舞台}{wu3tai2}[][HSK 3]
    \definition[个]{s.}{palco; plataforma elevada usada exclusivamente para apresentações artísticas, geralmente localizada na parte frontal de teatros e auditórios | palco; campo das atividades sociais}
  \synonymref{剧场}{ju4chang3}
  \synonymref{领域}{ling3yu4}
  \synonymref{平台}{ping2tai2}
  \synonymref{阵地}{zhen4di4}
  \antonymref{后台}{hou4tai2}
  \antonymref{幕后}{mu4hou4}
  \end{Phonetics}
\end{Entry}

\begin{Entry}{舞会}{14,6}{⾇,⼈}
  \begin{Phonetics}{舞会}{wu3hui4}
    \definition[场,个]{s.}{baile; dança; festa dançante; atividades em que muitas pessoas dançam juntas costumavam exigir trajes formais, mas agora são mais descontraídas e informais}
  \end{Phonetics}
\end{Entry}

\begin{Entry}{舞会舞}{14,6,14}{⾇,⼈,⾇}
  \begin{Phonetics}{舞会舞}{wu3hui4 wu3}
    \definition{s.}{baile}
  \end{Phonetics}
\end{Entry}

\begin{Entry}{舞抃}{14,7}{⾇,⼿}
  \begin{Phonetics}{舞抃}{wu3 bian4}
    \definition{s.}{Literário: bater palmas e dançar de alegria}
  \end{Phonetics}
\end{Entry}

\begin{Entry}{舞蹈}{14,17}{⾇,⾜}
  \begin{Phonetics}{舞蹈}{wu3dao3}[][HSK 6]
    \definition[段,支,场,个]{s.}{dança; uma forma de arte que usa movimentos rítmicos como principal meio de expressão, podendo expressar a vida, os pensamentos e os sentimentos das pessoas, geralmente acompanhada de música}
    \definition{v.}{dançar}
  \synonymref{跳舞}{tiao4//wu3}
  \end{Phonetics}
\end{Entry}

%%%%% EOF %%%%%

